Vision--language models (VLMs) combine a pretrained vision encoder with a large language model and achieve strong performance across a broad range of multimodal tasks. Yet they remain unreliable on tasks that require fine-grained, spatially grounded understanding (counting, referential grounding, and spatial reasoning), where predictions often reflect shortcut correlations and scene-level priors rather than evidence drawn from the queried region~\citep{chenspatial, sengupta2025can}.

A growing body of work locates part of this failure in how VLMs allocate visual attention. The attention from a generated token to the image patches can be read out as a \emph{saliency map}, and the degree to which this map overlaps the relevant object correlates with answer correctness~\citep{chenspatial}. In practice these maps are frequently mislocalized: attention disperses over non-informative regions, image borders, and background, and only a small subset of specialized heads attend in a grounding-relevant way~\citep{kangsee, kang2025your, kim2025interpreting}. Tellingly, \emph{shaping} the attention distribution (suppressing distractor regions) helps more than simply increasing the attention mass on targets~\citep{sengupta2025can}. Attention is thus both a faithful diagnostic of grounding and a promising intervention point.

Existing remedies stop short of supervising this signal directly. Inference-time methods reallocate or mask attention with hand-designed heuristics but never change what the model has learned. Representation-alignment objectives, which regularize a VLM's hidden states toward a pretrained vision encoder~\citep{yurepresentation, yoon2025visual}, do change the model, but their target is a \emph{static}, query-agnostic feature map: it cannot express that a token should attend \emph{here} for \emph{this} prompt. What is missing is a training-time objective that supervises the token-conditioned attention itself against where the evidence actually is.

We close this gap with the \textbf{Saliency Alignment Loss}: an auxiliary objective that, during fine-tuning, pulls each generated token's saliency map toward the image region its referent occupies. Using a densely grounded captioning corpus, COCONut-PanCap~\citep{deng2025coconut}, we attach a target distribution (uniform over the token's annotated mask) to every grounded token and minimize the divergence from the model's saliency through a Kullback--Leibler term added to the usual captioning loss. The objective introduces no new parameters and no architectural change, and because grounding benchmarks are overwhelmingly derived from COCO~\citep{han2024coco} (and therefore overlap our training images), we evaluate alignment \emph{intrinsically}, on held-out images, with three complementary attention-localization metrics.

Fine-tuning LLaVA-1.5-7B for only $200$ steps already reshapes its attention: all three localization metrics improve sharply (mean Attention Mass Ratio $4.6{\times}$) and the maps visibly concentrate on the objects a token names. We then ask \emph{where} this adaptation lives, and find, by parameter-drift analysis, that the language model, not the multimodal projector, absorbs the gradient, reaching a solution that is strikingly reproducible across runs and that preferentially moves a subset of attention heads, echoing the localization-head hypothesis~\citep{kang2025your}. Our most consequential finding, however, is what does \emph{not} change: despite the large intrinsic gains, downstream benchmarks barely move, saliency-map quality and task grounding turn out to be \emph{dissociated} (across configurations even anti-correlated with counting), and the objective withdraws attention from broad background ``stuff'' (snow, sky) the annotations never cover. We therefore present alignment as an effective but double-edged tool, and outline an instruction-tuned variant that ports the objective into the standard LLaVA-1.5 recipe to target the transfer gap.

\paragraph{Contributions.}
\begin{itemize}
    \item We propose the Saliency Alignment Loss, a parameter-free, architecture-agnostic training objective that supervises a VLM's token-conditioned attention with region-level grounding annotations (\Cref{sec:method}), together with an intrinsic, contamination-aware protocol (AMR, AP, NSS) for measuring attention localization on held-out data (\Cref{sec:metrics}).
    \item Our central empirical finding is a \emph{dissociation}: alignment improves the intrinsic metrics by large margins yet leaves downstream benchmarks essentially unchanged, and, across training configurations, attention quality is even anti-correlated with counting accuracy. Sharper attention is not, on its own, better grounding (\Cref{sec:results,sec:analysis}).
    \item We localize \emph{where} the model adapts through parameter-drift analysis: the language model rather than the projector absorbs the signal, identically with or without the projector trainable and reproducibly across runs; and we characterize a background-``stuff'' failure mode and root it in the annotation distribution (\Cref{sec:analysis}).
\end{itemize}
